\chapter{Судалгааны хэсэг}

\section{Компьютерын харааны тухай ойлголт, үүсэл хөгжил}
Компьютерын хараа (Computer Vision) нь дижитал зураг, дүрс бичлэгээс утга учир бүхий мэдээлэл автоматаар гарган авах, ойлгох, шийдвэр дэмжих зорилготой хиймэл оюуны үндсэн чиглэл юм \cite{real_time_detection}. Энэ салбарын хөгжлийг зөвхөн алгоритмын нарийвчлалаар бус, тооцооллын нөөц, өгөгдлийн хэмжээ, хэрэглээний орчны шаардлагатай нь хамтатган авч үзэх шаардлагатай. Ялангуяа жижиглэн худалдааны орчинд системийн үнэ цэнэ нь зөвхөн өндөр accuracy-гаар бус, бодит хугацаанд ажиллах чадвар, deployment өртөг, бага нөөцийн орчинд хэрэгжих боломжтойгоор хэмжигдэнэ.

Компьютерын харааны хөгжлийг дараах гурван үе шатанд авч үзэж болно.

\begin{enumerate}
\item \textbf{1960--2000 оны уламжлалт үе:} SIFT, HOG, SURF зэрэг гар аргаар тодорхойлсон feature-д тулгуурласан аргууд давамгайлж байсан. Эдгээр нь хяналттай орчинд үр дүнтэй боловч гэрэлтүүлэг, өнцөг, халхлалт өөрчлөгдөхөд тогтвортой биш байв.
\item \textbf{2012 оноос хойших гүн сургалтын үе:} AlexNet-ийн амжилтаас хойш CNN-д суурилсан дүрслэлийн суралцах арга хурдацтай хөгжиж, feature engineering-ийг feature learning-ээр орлуулсан \cite{alexnet}. Үүний үр дүнд ангилал, сегментаци, илрүүлэлтийн нарийвчлал эрс сайжирсан.
\item \textbf{Сүүлийн үеийн foundation-model чиг хандлага:} CLIP, SAM зэрэг foundation model-ууд дүрс-тайлбарын холбоо, ерөнхий сегментаци, zero-shot чадвараар салбарыг урагшлуулж байна \cite{clip,sam}. Гэвч ийм загваруудын санах ой, тооцоолол, inference latency өндөр тул мобайл ба хөнгөн серверийн орчинд шууд ашиглахад хүнд хэвээр байна. Иймээс энэхүү судалгаанд lightweight detector сонгох нь онолын болон инженерийн аль аль талаасаа үндэслэлтэй.
\end{enumerate}

Иймээс орчин үеийн компьютерын харааны судалгаа нэг талаас ерөнхий мэдлэгтэй том загвар, нөгөө талаас тодорхой домэйнд зориулсан хурдан, хөнгөн detector гэсэн хоёр чиглэлээр зэрэг хөгжиж байна. Энэхүү дипломын ажил нь хоёр дахь чиглэлд, өөрөөр хэлбэл retail shelf audit гэх тодорхой хэрэглээнд зориулсан practical detector сонгох нөхцөлд хамаарна.

\subsection{Конволюцийн мэдрэлийн сүлжээний математик үндэс}
CNN нь дүрслэлийн локал хэв шинжийг багаас их түвшинд шаталсан байдлаар суралцдаг тул объект илрүүлэлтийн суурь хэвээр байна \cite{cnn_lecun}. Конволюцийн үйлдэл нь оролтын зураг $I$ болон шүүлтүүр $K$-ийн хувьд дараах байдлаар тодорхойлогдоно:

\begin{equation}
    (I * K)(i, j) = \sum_{m} \sum_{n} I(i-m, j-n) \cdot K(m, n)
\end{equation}

Энд $(i, j)$ нь гаралтын пикселийн координат, $(m, n)$ нь шүүлтүүрийн индексүүдийг илэрхийлнэ. CNN-ийн давхаргууд нь зөвхөн feature extraction хийх биш, detection architecture-ийн дараагийн шатуудад ашиглагдах representation-ийг бэлтгэдэг.

\begin{enumerate}
\item \textbf{Конволюцийн давхарга:} Feature map үүсгэж, ирмэг, хэлбэр, texture зэрэг локал мэдээллийг ялгана. Гаралтын хэмжээ:
\begin{equation}
O = \frac{W - K + 2P}{S} + 1
\end{equation}
Энд $W$ нь оролтын хэмжээ, $K$ нь kernel, $P$ нь padding, $S$ нь stride юм.
\item \textbf{Идэвхжүүлэх функц:} ReLU нь шугаман бус хамаарлыг нэвтрүүлж, сургалтыг тогтвортой болгодог:
\begin{equation}
f(x) = \max(0, x)
\end{equation}
\item \textbf{Pooling давхарга:} Spatial хэмжээг багасгаж, локал өөрчлөлтөд тэсвэртэй болгоно:
\begin{equation}
y_{i,j} = \max_{(m,n) \in R_{i,j}} x_{m,n}
\end{equation}
\item \textbf{Бүрэн холбогдсон буюу prediction head:} Detection architecture-д ангилал ба байршлын таамгийг гаргах сүүлчийн шатны үндэс болдог. Softmax нь олон ангиллын магадлалыг дараах байдлаар тооцно:
\begin{equation}
\sigma(z)_i = \frac{e^{z_i}}{\sum_{j=1}^{K} e^{z_j}}
\end{equation}
\end{enumerate}

\begin{figure}[H]
\centering
\begin{tikzpicture}[
    node distance=0.6cm,
    box/.style={rectangle, draw, fill=blue!20, minimum width=1.2cm, minimum height=2cm, align=center, font=\scriptsize},
    smallbox/.style={rectangle, draw, fill=green!20, minimum width=0.8cm, minimum height=1.5cm, align=center, font=\scriptsize},
    arrow/.style={->, thick, >=stealth}
]

\node[box, fill=cyan!30, minimum height=2.5cm] (input) at (0, 0) {Input\\32x32x3};
\node[box, fill=blue!30, minimum height=2.2cm] (conv1) at (2, 0) {Conv\\28x28\\x32};
\node[smallbox, fill=orange!30, minimum height=1.8cm] (pool1) at (3.5, 0) {Pool\\14x14\\x32};
\node[box, fill=blue!30, minimum height=1.6cm] (conv2) at (5, 0) {Conv\\10x10\\x64};
\node[smallbox, fill=orange!30, minimum height=1.2cm] (pool2) at (6.5, 0) {Pool\\5x5\\x64};
\node[smallbox, fill=yellow!30, minimum height=0.8cm, minimum width=0.4cm] (flat) at (7.8, 0) {Flat\\1600};
\node[smallbox, fill=purple!30, minimum height=1cm, minimum width=0.5cm] (fc1) at (9, 0) {FC\\128};
\node[smallbox, fill=red!30, minimum height=0.6cm, minimum width=0.4cm] (fc2) at (10.2, 0) {FC\\8};

\draw[arrow] (input) -- (conv1);
\draw[arrow] (conv1) -- (pool1);
\draw[arrow] (pool1) -- (conv2);
\draw[arrow] (conv2) -- (pool2);
\draw[arrow] (pool2) -- (flat);
\draw[arrow] (flat) -- (fc1);
\draw[arrow] (fc1) -- (fc2);

\node[font=\scriptsize, below] at (1, -1.8) {Конволюци};
\node[font=\scriptsize, below] at (3.5, -1.8) {Pooling};
\node[font=\scriptsize, below] at (9.6, -1.8) {Prediction};

\end{tikzpicture}
\caption{CNN-ийн ерөнхий архитектурын бүтэц}
\label{fig:cnn_architecture}
\end{figure}

\section{Объект илрүүлэлтийн алгоритмуудын дэвшил}
Объект илрүүлэлт нь зурган дахь объектын ангилал болон байршлыг нэгэн зэрэг тодорхойлдог даалгавар бөгөөд accuracy-speed trade-off хамгийн мэдрэмтгий чиглэлүүдийн нэг юм. Энэ хөгжлийг two-stage ба one-stage ангиллаас гадна anchor-based-аас anchor-free руу шилжсэн өөрчлөлтөөр тайлбарлах нь илүү оновчтой.

\subsection{Хоёр алхмаар илрүүлдэг алгоритмууд}
Two-stage detector-ууд нь эхлээд candidate region санал болгож, дараа нь тэдгээрийг ангилж, bounding box-ийг нарийвчилдаг.

\begin{itemize}
\item \textbf{R-CNN:} Region proposal-уудыг тусад нь гаргаж, бүс тус бүрийг CNN-ээр үнэлдэг тул нарийвчлал сайтай боловч тооцооллын зардал маш өндөр байсан.
\item \textbf{Fast R-CNN:} Shared feature map ашигласнаар илүү үр ашигтай болсон ч region proposal шат тусдаа хэвээр үлдсэн.
\item \textbf{Faster R-CNN:} Region Proposal Network-ийг нэгтгэснээр end-to-end сургалтын боломж нээгдсэн \cite{faster_rcnn}. Гэсэн хэдий ч олон шаттай pipeline нь inference latency өндөр хэвээр байв.
\end{itemize}

Эдгээр судалгааны гол ач холбогдол нь өндөр нарийвчлалтай detection pipeline бий болгосонд оршино. Гэвч жижиглэн худалдааны мобайл аудитын нөхцөлд latency болон deployment simplicity нь илүү чухал тул two-stage загваруудын онолын давуу тал нь практик сонголт болж чаддаггүй.

\subsection{Нэг алхмаар илрүүлдэг алгоритмууд}
One-stage detector-ууд нь proposal generation ба classification/regression-ийг нэг урсгалд нэгтгэн, real-time хэрэглээнд илүү тохиромжтой.

\begin{itemize}
\item \textbf{SSD:} Multi-scale feature map ашиглан янз бүрийн хэмжээтэй объект илрүүлэхийг зорьсон \cite{ssd,fpn}. Хурд сайжирсан боловч densely packed shelf орчинд жижиг объектын ялгарал сул байж болдог.
\item \textbf{YOLO:} Detection-ийг нэг pass-аар хийснээр object detection-ийг real-time хэрэглээнд ойртуулсан \cite{yolo_original}. Дараагийн хөгжлийн шугам нь энэ хурдыг алдагдуулахгүйгээр нарийвчлалыг өсгөхөд чиглэсэн.
\item \textbf{RetinaNet:} Focal Loss ашиглан class imbalance-ийн асуудлыг бууруулж, one-stage detector-уудын нарийвчлалыг өсгөсөн \cite{focal_loss}.
\end{itemize}

\subsection{Anchor-based ба anchor-free хувьсал}
Объект илрүүлэлтийн архитектурын өөр нэг чухал хөгжлийн шугам нь anchor-based-аас anchor-free руу шилжсэн явдал юм. Anchor-based загварууд нь объектын боломжит хэмжээ, aspect ratio-ийг урьдчилан тодорхойлсон anchor box-уудаар төлөөлүүлдэг. Энэ арга нь олон жилийн турш үр дүнтэй байсан боловч дараах хүндрэл дагуулдаг: (i) anchor хэмжээ, харьцааг гар аргаар тааруулах шаардлагатай, (ii) dense prediction үед negative sample ихсэж class imbalance нэмэгддэг, (iii) шинэ домэйнд дахин тохируулга хийх инженерийн зардал өндөр байдаг.

Anchor-free загварууд нь объектын төв, keypoint эсвэл boundary prediction-д тулгуурлан эдгээр хүндрэлүүдийг багасгахыг зорьдог. Энэ чиг хандлага нь one-stage detector-уудыг илүү цэвэр optimization-той болгож, deployment-oriented хэрэглээнд давуу тал өгсөн. Ge нарын YOLOX ажил нь anchor-free head-ийг YOLO гэр бүлийн хүрээнд бодит гүйцэтгэлийн сайжруулалттай холбон харуулсан төлөөлөх жишээ бөгөөд \cite{yolox} энэхүү хувьсал нь YOLO-ийн сүүлийн үеийн хувилбаруудад шууд нөлөөлсөн. Иймээс YOLOv8-ийг энэ чиглэлийн логик үргэлжлэл гэж үзэж болно.

\section{YOLO алгоритмын хувилбаруудын хөгжил}
YOLO алгоритмын хөгжлийг зөвхөн хувилбаруудын жагсаалт бус, ``хурд, нарийвчлал, хэрэгжүүлэхэд хялбар байдал'' гэсэн гурван шалгуурын тасралтгүй тэнцвэржүүлэлт гэж үзэх нь зүйтэй.

\begin{longtable}{|l|c|c|c|p{5.2cm}|}
\caption{YOLO хувилбаруудын харьцуулалт}
\label{tab:yolo_versions}\\
\hline
\textbf{Хувилбар} & \textbf{Он} & \textbf{mAP (\%)} & \textbf{FPS} & \textbf{Онцлог шинж} \\
\hline
\endfirsthead
\hline
\textbf{Хувилбар} & \textbf{Он} & \textbf{mAP (\%)} & \textbf{FPS} & \textbf{Онцлог шинж} \\
\hline
\endhead
\hline
\endfoot
YOLOv1 & 2015 & 63.4 & 45 & Detection-ийг нэг шатанд оруулсан анхны практик хувилбар \cite{yolo_original} \\
\hline
YOLOv2 & 2016 & 76.8 & 67 & Batch normalization, anchor box ашиглалт сайжирсан \cite{yolov2} \\
\hline
YOLOv3 & 2018 & 82.3 & 35 & Multi-scale detection, feature pyramid сэтгэлгээ хүчтэй болсон \cite{yolov3} \\
\hline
YOLOv4 & 2020 & 84.5 & 62 & CSPDarknet53, PANet ашиглаж speed-accuracy тэнцвэржүүлсэн \cite{yolov4} \\
\hline
YOLOv5 & 2020 & 86.2 & 140 & PyTorch экосистемд өргөн хэрэглээнд орсон, инженерийн хувьд боловсорсон \\
\hline
YOLOv8 & 2023 & 88.9 & 280 & Anchor-free head, C2f модуль, Ultralytics-ийн нэгдсэн training/inference pipeline \cite{yolov8_ultralytics} \\
\hline
\end{longtable}

Энэ судалгаанд YOLOv8-ийг сонгосон шалтгаан нь ``хамгийн шинэ'' байсан учраас биш, харин судалгааны ажлын хэрэгжилтийн хугацаанд хамгийн тэнцвэртэй сонголт байсан явдал юм. Энэхүү ажлын датасет бүрдүүлэлт, туршилтын pipeline тогтож байсан үед YOLOv8 нь:

\begin{enumerate}
\item Ultralytics экосистем дотор баримтжуулалт, жишээ код, training workflow-ийн хувьд хамгийн боловсорсон,
\item anchor-free architecture руу шилжсэн тул шинэ домэйнд anchor tuning хийх шаардлага багатай,
\item жижиг хэмжээний `n` хувилбар нь мобайл-бага нөөцийн хэрэглээнд тохиромжтой,
\item судалгааны reproducibility-г хангахуйц өргөн хэрэглээтэй байсан.
\end{enumerate}

YOLOv9, YOLOv10 зэрэг дараагийн хувилбарууд нь илүү сүүлийн үеийн дэвшил авчирсан ч энэхүү судалгааны temporal scope дотор бүх baseline, implementation stack-ийг тэдгээр рүү шилжүүлэн дахин баталгаажуулах нь ажлын үндсэн зорилтоос хазайх эрсдэлтэй байв. Иймээс энэ ажилд YOLOv8-ийг implementation maturity болон domain adaptation-ийн хувьд хамгийн зохистой суурь гэж үзсэн.

YOLOv8-ийн гол инженерийн давуу талууд нь дараах байдалтай.

\begin{enumerate}
\item \textbf{Anchor-free detection:} Domain-specific anchor design шаардлагагүй тул сургалт, туршилтын мөчлөгийг богиносгосон.
\item \textbf{C2f module:} Параметрийн үр ашгийг сайжруулж, compact model-уудад feature flow-ийг тогтвортой болгосон.
\item \textbf{Decoupled head:} Classification ба localization-ийг тусгаарласан нь shelf object detection зэрэг нарийн ялгаралтай орчинд ашигтай.
\item \textbf{Нэгдсэн augmentation болон training ecosystem:} Mosaic, MixUp зэрэг augmentation-ууд нь жижиг датасетийн нөхцөлд generalization сайжруулахад дэмжлэг болсон \cite{data_augmentation}.
\end{enumerate}

\begin{figure}[H]
\centering
\begin{tikzpicture}[
    node distance=1.5cm,
    box/.style={rectangle, draw, fill=blue!20, minimum width=1.8cm, minimum height=0.8cm, align=center, font=\footnotesize},
    arrow/.style={->, thick}
]
\draw[thick, ->] (0,0) -- (13,0) node[right] {Он};
\foreach \x/\year in {1/2015, 3/2016, 5/2018, 7/2020, 9/2021, 11/2023} {
    \draw (\x, 0.1) -- (\x, -0.1) node[below, font=\footnotesize] {\year};
}
\node[box, fill=red!30] at (1, 1.2) {YOLOv1};
\node[box, fill=orange!30] at (3, 1.2) {YOLOv2};
\node[box, fill=yellow!30] at (5, 1.2) {YOLOv3};
\node[box, fill=green!30] at (7, 1.2) {YOLOv4};
\node[box, fill=cyan!30] at (9, 1.2) {YOLOv5};
\node[box, fill=purple!30] at (11, 1.2) {YOLOv8};
\node[align=center, font=\scriptsize] at (1, 2.2) {Single-stage\\detection};
\node[align=center, font=\scriptsize] at (7, 2.2) {Speed-accuracy\\refinement};
\node[align=center, font=\scriptsize] at (11, 2.2) {Anchor-free\\maturity};
\end{tikzpicture}
\caption{YOLO алгоритмын хувилбаруудын хөгжлийн цаг хугацааны диаграмм}
\label{fig:yolo_timeline}
\end{figure}

\section{Жижиглэн худалдааны салбар дахь ухаалаг мониторингийн системүүд}
Энэ дэд хэсэг нь бүрэн systematic literature review биш боловч retail shelf monitoring сэдэвт төлөөлөх судалгаа, үйлдвэрлэлийн шийдлүүдийг бүтэцтэйгээр авч үзнэ. Тухайлбал 2017--2024 оны хоорондын төлөөлөх академик өгүүлэл, бодит хэрэгжүүлэлтийг хамруулж, дараах асуултад төвлөрөв: (i) систем ямар орчинд ажилладаг вэ, (ii) ямар өгөгдөл ашигладаг вэ, (iii) эцсийн хэрэглээнд хэр хэмжээнд end-to-end хүрдэг вэ?

Santra болон Mukherjee нарын retail product recognition-ийн тойм нь энэ чиглэлийн судалгааг өргөн цар хүрээнд нэгтгэж, CNN-д суурилсан аргууд уламжлалт аргуудаас тогтвортой илүү үр дүнтэй байгааг харуулсан \cite{retail_cv_survey}. Гэвч ийм тоймуудын түгээмэл сул тал нь бодит өртөг, deployment complexity, локал домэйнд дасан зохицох чадварыг харьцангуй бага авч үздэгт оршино.

Goldman нарын SKU-110K датасет нь нягт өрөгдсөн дэлгүүрийн барааны илрүүлэлтэд зориулсан чухал benchmark болсон \cite{sku110k}. Гэвч ийм benchmark дээрх сайн үр дүн нь өөр бүс нутгийн бараа, савлагаа, гэрэлтүүлэгт шууд шилжинэ гэсэн үг биш. Retail domain-д dataset shift, packaging shift, camera-angle shift түгээмэл тохиолддог.

\subsection{Үйлдвэрлэлийн бодит шийдлүүд}
\textbf{Amazon Go} нь өндөр нягтралтай камер, мэдрэгч, тусгай дэд бүтэц дээр тулгуурласан бүрэн автомат шийдэл юм \cite{amazon_go}. Frictionless retail боломжийг баталсан ч өртөг, дэд бүтцийн төвөгшил нь жижиг, дунд байгууллагад шууд хэрэгжүүлэхэд саад болдог.

\textbf{Trax Retail} нь ухаалаг төхөөрөмжөөр зураг цуглуулж, клоуд дээр боловсруулан retail execution audit хийдэг арилжааны шийдэл \cite{trax_retail}. Энэ нь мобайл суурьтай audit workflow бодит хэрэглээнд хэрэгжих боломжтойг харуулдаг. Гэвч хаалттай эх, лицензийн өртөг, алгоритмын ил тод бус байдал нь академик болон локал customization-д хязгаар үүсгэнэ.

\textbf{Bossa Nova Robotics} нь робот платформоор дэлгүүрийн эгнээг автомат скан хийх хандлагыг төлөөлдөг. Энэ чиглэл нь хүний оролцоог багасгадаг боловч нарийн гарцтай дэлгүүр, бага төсөвтэй зах зээлд operational fit сул байдаг.

Дүгнэж хэлбэл, үйлдвэрлэлийн шийдлүүдийн ихэнх нь өндөр автоматжуулалттай боловч өндөр өртөгтэй, эсвэл уян хатан боловч хаалттай байдаг. Иймээс ``low-cost, locally trainable, end-to-end'' шийдлийн орон зай нээлттэй хэвээр байна.

\section{Ижил төстэй системүүдийн харьцуулсан шинжилгээ}
Related work-ийг зөв байрлуулахын тулд зөвхөн системүүдийг нэрлэх бус, тэдгээрийн өгөгдөл, архитектур, хэрэгжүүлэх цар хүрээг харьцуулах шаардлагатай. Доорх хүснэгт нь үйлдвэрлэлийн болон академик төлөөлөх шийдлүүдийг энэхүү дипломын ажлын зорилготой харьцуулан үзүүлэв.

\begingroup
\footnotesize
\setlength{\tabcolsep}{3pt}
\renewcommand{\arraystretch}{1.06}
\begin{longtable}{|P{0.15}|P{0.11}|P{0.14}|P{0.13}|P{0.09}|P{0.12}|P{0.18}|}
\caption{Жижиглэн худалдааны ухаалаг системүүдийн харьцуулсан шинжилгээ}
\label{tab:system_comparison}\\
\hline
\textbf{Систем / Судалгаа} & \textbf{Өгөгдөл} & \textbf{Орчин} & \textbf{Орон нутгийн тохируулга} & \textbf{End-to-end} & \textbf{Төхөөрөмж} & \textbf{Гол хязгаарлалт} \\
\hline
\endfirsthead
\hline
\textbf{Систем / Судалгаа} & \textbf{Өгөгдөл} & \textbf{Орчин} & \textbf{Орон нутгийн тохируулга} & \textbf{End-to-end} & \textbf{Төхөөрөмж} & \textbf{Гол хязгаарлалт} \\
\hline
\endhead
\hline
\endfoot
Amazon Go & Нээлттэй биш & Ухаалаг дэлгүүр & Үгүй & Тийм & Олон камер, мэдрэгч & Суурилуулалтын өртөг маш өндөр; ердийн дэлгүүрт нэвтрүүлэхэд төвөгтэй \\
\hline
Trax Retail & Нээлттэй биш & Клоуд audit & \shortstack{Хязгаар\\лагдмал} & Хэсэгчлэн & Ухаалаг утас / таблет & Хаалттай экосистем; лицензийн өртөг өндөр \\
\hline
Bossa Nova & Нээлттэй биш & Робот скан & \shortstack{Хязгаар\\лагдмал} & Тийм & Робот платформ & Ашиглалтын төвөгшил өндөр; жижиг орчинд тохиромж муутай \\
\hline
Zhang et al. \cite{shelf_monitoring} & Дунд хэмжээний зураг & Shelf monitoring & \shortstack{Хязгаар\\лагдмал} & Хэсэгчлэн & Камер + сервер & Домэйнд тохируулах арга, мобайл хэрэгжилт тодорхой бус \\
\hline
Tonioni et al. \cite{planogram_compliance} & Planogram өгөгдөл & Shelf compliance & \shortstack{Хязгаар\\лагдмал} & Үгүй & Камер + graph matching & Илрүүлэлтээс илүү planogram matching-д төвлөрсөн \\
\hline
Энэ ажил & Орон нутгийн өгөгдөл & \shortstack{Монголын жижиглэн\\худалдааны үйлчилгээний цэг} & Тийм & Тийм & Ухаалаг утас + хөнгөн сервер & Ангиллын тоо, өгөгдлийн хамрах хүрээ одоогоор \shortstack{хязгаар\\лагдмал} \\
\hline
\end{longtable}
\endgroup

Энэ хүснэгтээс хоёр зүйл тод харагдана. Нэгдүгээрт, академик судалгаануудын ихэнх нь detection accuracy эсвэл planogram logic-д төвлөрдөг боловч бодит ажлын урсгалтай нэгтгэсэн end-to-end систем ховор. Хоёрдугаарт, үйлдвэрлэлийн шийдлүүд бодит хэрэглээнд боловсорсон ч нээлттэй, хямд өртөгтэй, локал бараанд дасан зохицох боломж бага. Энэхүү дипломын ажил нь эдгээр хоёр зайны завсарт орших practical gap-ийг нөхөх зорилготой.

\section{Бидний системийн давуу тал ба байр суурь}
Энэхүү системийн давуу талыг тусгаарлан магтах байдлаар бус, өмнөх ажлуудтай харьцуулан үнэлэх нь илүү академик зөв хандлага юм. Энэ үүднээс дараах байршуулалтыг хийж болно.

\begin{enumerate}
\item \textbf{Өртөг ба хэрэгжүүлэх төвөгшлийн хувьд:} Amazon Go болон робот суурьтай шийдлүүдтэй харьцуулахад ухаалаг утас + хөнгөн серверийн зохиомж нь дэд бүтцийн босгыг мэдэгдэхүйц бууруулдаг.
\item \textbf{Нээлттэй, дахин ашиглах боломжийн хувьд:} Trax зэрэг хаалттай шийдлүүдээс ялгаатай нь нээлттэй эхийн YOLO экосистем дээр тулгуурласан тул сургалт, fine-tuning, өргөтгөлийг судалгааны орчинд хянах боломжтой.
\item \textbf{Домэйн тохируулгын хувьд:} Олон улсын benchmark-д суурилсан ажлуудаас ялгаатай нь энэхүү систем локал retail shelf өгөгдөлд тулгуурласан тул domain shift-ийг бууруулах практик ач холбогдолтой.
\item \textbf{End-to-end audit workflow-ийн хувьд:} Зарим related work зөвхөн detection эсвэл compliance-checking төвтэй байдаг бол энэ ажилд зураг авах, илрүүлэлт, аудитын тооцоолол, үр дүн харуулах бүтэн урсгал холбогдсон.
\item \textbf{Судалгаа ба хэрэгжүүлэлтийн завсрын байрлалын хувьд:} Энэ ажил нь зөвхөн benchmark score хөөсөн бус, харин бодит хэрэглээнд шилжих боломжтой архитектур санал болгосноор хэрэгжүүлэлтийн цоорхойг нөхөж байна.
\end{enumerate}

\section{Судалгааны цоорхой ба энэхүү ажлын байрлал}
Дээрх судалгаануудыг нэгтгэн дүгнэвэл дараах research gap тодорхойлогдож байна.

\begin{enumerate}
\item Retail shelf monitoring-ийн академик судалгаанууд өндөр нарийвчлалтай detection эсвэл planogram verification-д төвлөрөх хандлагатай боловч мобайл хэрэглэгчийн урсгалтай end-to-end аудитын системийг бүрэн авч үзэх нь ховор.
\item Үйлдвэрлэлийн бодит шийдлүүд хэрэгжүүлэлтийн хувьд боловсорсон ч хаалттай, өндөр өртөгтэй, тухайн орны бараа бүтээгдэхүүнд тохируулан дахин сургах боломж бага байдаг.
\item Нээлттэй, lightweight detector дээр суурилсан, орон нутгийн датасеттэй, жижиглэн худалдааны бодит audit workflow-тай уялдсан системийн судалгаа харьцангуй цөөн байна.
\item Олон улсын benchmark датасет дээрх сайн үр дүн нь Монголын retail domain-д шууд шилжинэ гэсэн баталгаа биш бөгөөд локал adaptation нь өөрөө судалгааны үнэ цэнтэй асуудал хэвээр байна.
\end{enumerate}

Иймээс энэхүү дипломын ажил нь YOLOv8 суурьтай lightweight detector-ийг орон нутгийн retail shelf өгөгдөлд тохируулж, гар утаснаас зураг авч, сервер талд илрүүлэлт хийж, аудитын дүгнэлт автоматаар гаргах end-to-end систем боловсруулах замаар дээрх цоорхойг нөхөхийг зорьж байна. Өөрөөр хэлбэл, энэ ажил нь ``хамгийн том загвар'' бүтээхийг бус, ``хамгийн бодитой хэрэгжих боломжтой локал шийдэл'' боловсруулахыг судалгааны үндсэн байр сууриа болгон авч байна.
